\section{Technical Approach}
We engineer a ZXW-native intermediate representation that embeds the checkpoints advocated in ``Mind the gaps'' directly into the compilation workflow.

\subsection{ZXW Canonicalisation Layer}
\begin{itemize}
  \item Extend QUIZX with W-spider rewrite schemas (absorption, fusion, controlled phase propagation) whose soundness is machine-checked and logged, turning heuristic rewrites into verifiable artefacts.
  \item Introduce scalable diagram notation that compresses repeated ZXW motifs while encoding provenance metadata for downstream simulators and code designers.
  \item Instrument the layer with error-mitigation-aware metrics (node counts, depth, stochastic noise channels) to quantify impact and inform when active error detection should take over.
\end{itemize}

\subsection{Hybrid Feedback and Optimisation}
\begin{itemize}
  \item Embed classical optimisation heuristics (flow, gflow, stabiliser checks) that target code-compatible layouts and surface when transitions to QEC-ready circuits are warranted.
  \item Maintain a certificate log that records rule applications, referenced lemmas, and hardware-specific constraints, enabling reproducibility and easier validation against classical co-simulators.
  \item Provide hooks for numerical routines that evaluate expectation values on tractable instances, creating ``proof pockets'' that substantiate claimed gains before full-scale experiments.
\end{itemize}

\subsection{Gate-Set and Code-Aware Extraction}
\begin{itemize}
  \item Translate canonical diagrams into circuits respecting restricted gate alphabets and lattice layouts relevant to superconducting, ion-trap, and neutral-atom architectures.
  \item Generate diagnostic reports that bundle depth, two-qubit counts, logical-to-physical mappings, and mitigation strategies, supporting the escalation from mitigation to detection and correction.
  \item Align extraction routines with benchmarking suites and error-budget models so that claims of simulation advantage remain auditable.
\end{itemize}
